\section{Solvable quintics and cyclotomic resolvents}\label{sec:solvable}

Abel--Ruffini leaves a large and interesting solvable class.  In this section
all displayed root identities are proved by finite expansion in a cyclotomic
extension.  The point is not that the resulting formulas are small, but that
their correctness is decidable by exact algebra.

\subsection{Reducible equations}

A reducible quintic is solved factor by factor.  For example,
\begin{equation}\label{eq:reducible-example}
 X^5-X^4-X+1=(X^2+1)(X+1)(X-1)^2.
\end{equation}
This identity is immediate by multiplication and is an \ipra{} theorem after
integer polynomial normalization.  Its roots use only the checked quadratic
formula and rational roots.  Such examples are not exceptions to
Theorem~\ref{thm:abel-ruffini}: that theorem concerns a uniform formula for the
generic quintic.

\subsection{The de Moivre quintic}

Let \(\omega\) satisfy~\eqref{eq:cyclo5}, and let \(u,v\) lie in a commutative
field containing \(\omega\).  Define
\[
 r_j=\omega^j u+\omega^{-j}v,
 \qquad 0\le j<5.
\]

\begin{theorem}[de Moivre product identity]\label{thm:de-moivre-product}
One has
\begin{equation}\label{eq:de-moivre-product}
 \prod_{j=0}^{4}(X-r_j)
 =X^5-5uvX^3+5u^2v^2X-(u^5+v^5).
\end{equation}
This identity is valid in the finite quotient
\(\Z[\omega,u,v,X]/(1+\omega+\cdots+\omega^4)\) and is therefore directly
formalizable in \ipra.
\end{theorem}

\begin{proof}
Let \(s_m=\sum_jr_j^m\).  Since
\(\sum_{j=0}^4\omega^{kj}=0\) unless \(5\mid k\), expansion gives
\[
 s_1=s_3=0,
 \qquad s_2=10uv,
 \qquad s_4=30u^2v^2.
\]
Newton's identities then yield
\[
 e_1=0,
 \quad e_2=-5uv,
 \quad e_3=0,
 \quad e_4=5u^2v^2.
\]
Finally,
\[
 e_5=\prod_jr_j
 =\prod_j(v+u\omega^{2j})=u^5+v^5,
\]
because multiplication by two permutes the fifth roots of unity and
\(\prod_{j=0}^4(Z+\omega^j)=Z^5+1\).  Substitution into
\(\prod(X-r_j)=X^5-e_1X^4+e_2X^3-e_3X^2+e_4X-e_5\) gives the result.
\end{proof}

Take \(uv=-a\) and \(u^5+v^5=-b\).  Then~\eqref{eq:de-moivre-product} becomes
\begin{equation}\label{eq:de-moivre-quintic}
 X^5+5aX^3+5a^2X+b.
\end{equation}
If \(Y=u^5\), then \(v=-a/u\) and
\[
 v^5=-\frac{a^5}{Y}.
\]
The equation \(u^5+v^5=-b\) is therefore equivalent to
\begin{equation}\label{eq:de-moivre-aux}
 Y^2+bY-a^5=0.
\end{equation}

\begin{corollary}[Radical solution of the de Moivre family]
\label{cor:de-moivre-roots}
Let \(Y\) be either root of~\eqref{eq:de-moivre-aux}, choose \(u^5=Y\), and
put \(v=-a/u\).  Then the five roots of~\eqref{eq:de-moivre-quintic} are
\begin{equation}\label{eq:de-moivre-roots}
 X_j=\omega^ju-\frac{a}{\omega^ju},
 \qquad 0\le j<5.
\end{equation}
Every defining relation and the final substitution are finite radical-tower
certificates.
\end{corollary}

The auxiliary equation is quadratic; the remaining operations are fifth roots
and fifth roots of unity.  Thus this family is solved by ordinary radicals in
the strict sense.

\subsection{The four-resolvent identity printed by MathWorld}

Let \(u_1,u_2,u_3,u_4\) be elements of a commutative cyclotomic extension and
put
\begin{equation}\label{eq:four-resolvent-roots}
 r_j=\omega^ju_1+\omega^{2j}u_2+\omega^{3j}u_3+\omega^{4j}u_4.
\end{equation}
Define
\begin{align*}
 U&=u_1u_4+u_2u_3,\\
 V&=u_1u_2^2+u_2u_4^2+u_3u_1^2+u_4u_3^2,\\
 W&=u_1^2u_4^2+u_2^2u_3^2-u_1^3u_2-u_2^3u_4-u_3^3u_1-u_4^3u_3
     -u_1u_2u_3u_4,\\
 X_0&=u_1^3u_3u_4+u_2^3u_1u_3+u_3^3u_2u_4+u_4^3u_1u_2,\\
 Y_0&=u_1u_3^2u_4^2+u_2u_1^2u_3^2+u_3u_2^2u_4^2+u_4u_1^2u_2^2,\\
 Z&=u_1^5+u_2^5+u_3^5+u_4^5.
\end{align*}

\begin{theorem}[Four-resolvent product]\label{thm:four-resolvent}
The five values~\eqref{eq:four-resolvent-roots} are the roots of
\begin{equation}\label{eq:four-resolvent-product}
 T^5-5UT^3-5VT^2+5WT+5(X_0-Y_0)-Z.
\end{equation}
\end{theorem}

\begin{proof}
Expand \(\prod_{j=0}^4(T-r_j)\), reduce every power of \(\omega\) with
\(1+\omega+\cdots+\omega^4=0\), and collect monomials in the \(u_i\).  The
coefficients are exactly those displayed above.  This finite calculation is
implemented independently by the archive's checker; its soundness in \ipra{}
is Lemma~\ref{lem:poly-normalize} in the cyclotomic quotient.  A human
alternative is to compute the first five power sums of the \(r_j\) and apply
Newton's identities, using the root-of-unity filter as in
Theorem~\ref{thm:de-moivre-product}.
\end{proof}

\begin{corollary}[Conditional root formula for \(T^5+aT+b\)]
\label{cor:four-resolvent-specialization}
If
\begin{equation}\label{eq:four-resolvent-constraints}
 U=0,
 \qquad V=0,
 \qquad W=\frac a5,
 \qquad 5(X_0-Y_0)-Z=b,
\end{equation}
then every \(r_j\) in~\eqref{eq:four-resolvent-roots} satisfies
\(r_j^5+ar_j+b=0\).
\end{corollary}

This proves the algebraic core of MathWorld's formulas (43)--(55).  The long
expressions there for the \(u_i\) are a way of manufacturing a solution of
\eqref{eq:four-resolvent-constraints}.  Their correctness can be checked by
adjoining symbols for the indicated square and fifth roots, clearing
nonzero denominators, and normalizing the four identities.  That is a finite
certificate problem; the present packet proves the uniform product identity
and specialization, while leaving the page's full parameter substitution as a
separate generated certificate rather than reproducing several pages of
radicals.

\subsection{Why every solvable irreducible quintic has resolvents of this kind}

The affine-group criterion gives a conceptual construction.  Label the roots
\(x_0,\ldots,x_4\) so that the normal translation subgroup \(C_5\) acts by
\(x_j\mapsto x_{j+1}\).  Define Lagrange resolvents
\begin{equation}\label{eq:lagrange-resolvents}
 R_k=\sum_{j=0}^4\omega^{-kj}x_j,
 \qquad 0\le k<5.
\end{equation}
Under the generating 5-cycle,
\[
 R_k\longmapsto\omega^kR_k.
\]
Therefore \(R_k^5\) is fixed by the normal \(C_5\).  The quotient of the
Galois group by \(C_5\) is a subgroup of \(C_4\), so each \(R_k^5\) lies in an
extension of degree at most four over the base field and can be obtained using
the checked formulas through Ferrari.  Once compatible fifth roots \(R_k\)
are chosen, Fourier inversion gives
\begin{equation}\label{eq:resolvent-inversion}
 x_j=\frac15\sum_{k=0}^4\omega^{kj}R_k.
\end{equation}

\begin{theorem}[Constructive resolvent solution]
\label{thm:solvable-quintic-resolvents}
Given an irreducible quintic together with a finite Galois certificate embedding
its group into \(\operatorname{AGL}(1,5)\), the construction
\eqref{eq:lagrange-resolvents}--\eqref{eq:resolvent-inversion} yields a finite
radical-tower certificate for all five roots.  Every compatibility choice is
verified by substitution rather than by an unbounded search.
\end{theorem}

\begin{proof}
The normal \(C_5\) and quotient action are read from the supplied permutation
table.  Elementary symmetric polynomials in the \(R_k^5\) are fixed by the
whole group and can be computed as rational functions of the original
coefficients by finite orbit sums.  The quotient degree is at most four, so the
checked quadratic, Cardano and Ferrari identities factor the corresponding
resolvent polynomial.  Adjoin one fifth root and obtain the remaining
compatible roots by the recorded group action.  Finally verify
\eqref{eq:resolvent-inversion} and the original quintic by cyclotomic
normalization.  Each stage is a bounded calculation on finite codes.
\end{proof}

This theorem explains the source page's statement that for prime degree five a
root may be written as a sum of fifth roots of the roots of a quartic auxiliary
equation, with compatibility relations among those fifth roots.

\subsection{Explicit radical examples}

The following examples illustrate how short the final verification can be even
when discovering the expression was difficult.

\begin{example}[A single fifth radical]
Let \(t^5=2\) and
\[
 x=1+t-t^2+t^3-t^4.
\]
Then
\begin{equation}\label{eq:explicit-example-one}
 x^5-5x^4+30x^3-50x^2+55x-21=0.
\end{equation}
Indeed, substituting the displayed \(x\) makes the left side a multiple of
\(t^5-2\).  Polynomial division gives an exact quotient of degree 15, recorded
by the executable certificate.
\end{example}

\begin{example}[A Paxton Young example]
With the same \(t^5=2\), put
\[
 x=t-t^2+t^3-t^4.
\]
Then
\begin{equation}\label{eq:young-example}
 x^5+20x^3+20x^2+30x+10=0.
\end{equation}
Again the substituted polynomial has zero remainder modulo \(t^5-2\).
\end{example}

\begin{example}[An eleventh-root-of-unity quintic]
Let \(z^{11}=1\), \(z\ne1\), and \(x=z+z^{-1}\).  Then
\begin{equation}\label{eq:cos11-example}
 x^5+x^4-4x^3-3x^2+3x+1=0.
\end{equation}
To verify this, multiply the left side after substituting \(x=z+z^{-1}\) by
\(z^5\) and reduce modulo
\(1+z+z^2+\cdots+z^{10}\).  The remainder is zero.  Under the usual complex
embedding, these roots are \(2\cos(2k\pi/11)\).
\end{example}

All three examples are exact quotient-ring computations and require no appeal
to numerical approximation.

\subsection{Fifth roots of unity by square roots}

The de Moivre formulas themselves use \(\omega\), but a primitive fifth root of
unity is obtainable by quadratic radicals.  Let \(s^2=5\), choose
\(\alpha,\beta\in\{-1,1\}\), let
\(t^2=-10-2\beta s\), and put
\begin{equation}\label{eq:fifth-root-radical}
 \omega_{\alpha,\beta}
 =\frac{\alpha t+\beta s-1}{4}.
\end{equation}
Clearing the denominator and reducing successively by
\(s^2-5\) and \(t^2+10+2\beta s\) proves
\[
 1+\omega_{\alpha,\beta}+\omega_{\alpha,\beta}^2
 +\omega_{\alpha,\beta}^3+\omega_{\alpha,\beta}^4=0.
\]
The four sign pairs give the four primitive roots; inequality of their real and
imaginary components is decidable in the corresponding ordered quadratic
subfields.  Thus every cyclotomic constant used above is itself represented by
a short radical certificate.

\subsection{Parameterized solvable trinomials}

The source pages quote the Runge--Cadenhad--Young and Spearman--Williams
parameterizations of solvable irreducible trinomials \(X^5+aX+b\).  Their
logical content separates naturally:

\begin{enumerate}
\item The displayed rational formulas for \(a,b\) are primitive-recursive
      maps of the rational parameters.
\item The displayed radical expressions for the \(u_i\), once supplied, are
      checked by Corollary~\ref{cor:four-resolvent-specialization} after
      clearing their square- and fifth-root relations.
\item The ``if and only if'' classification is a Galois theorem.  By
      Theorem~\ref{thm:AGL-criterion}, it reduces to constructing rational
      invariants that parameterize the embeddings of the particular trinomial
      Galois group into \(\operatorname{AGL}(1,5)\).  This is finite invariant
      theory, but its published formulas are lengthy.
\end{enumerate}

Accordingly, this packet proves the finite group criterion and the universal
root verifier, and treats the exact necessity/sufficiency substitution for the
published parameterization as a generated certificate target in the Lean
roadmap.  This is stronger than accepting the formulas numerically, but more
honest than claiming a handwritten formal derivation that has not been
expanded.

\subsection{Casus irreducibilis}

The source page notes that some solvable quintics have five real roots but any
radical formula passes through nonreal intermediate values.  The field-theoretic
mechanism is the same as for the cubic.  If an irreducible, totally real
extension were generated through a tower in which every radical and every
intermediate embedding remained real, then any odd-prime radical step would
have only one real conjugate and could not contribute to the normal totally
real closure.  After eliminating such steps, only quadratic extensions remain,
so the degree is a power of two.  An irreducible quintic has degree five, not a
power of two; hence a radical solution of a five-real-root solvable quintic must
use complex intermediates.

For \ipra{} this argument is formulated for coded finite ordered fields and
embedding tables: at an odd-prime radical step the other roots are
\(\zeta^j\alpha\), and the table certifies that they are nonreal; a tower that
keeps every conjugate real can therefore have only index two steps.  The full
ordered-field tower module is included as a formalization target rather than a
completed Lean file, so the audit labels this global theorem \(\Outline\), even
though every individual finite tower can already be checked by the certificate
interface.
